\documentclass[a4paper,12pt,twoside]{memoir}

% Castellano
\usepackage[spanish,es-tabla]{babel}
\selectlanguage{spanish}
\usepackage[utf8]{inputenc}
\usepackage[T1]{fontenc}
\usepackage{lmodern} % Scalable font
\usepackage{microtype}
\usepackage{placeins}
\usepackage{url} % Mejora el tratamiento de URL, p. ej. en saltos de línea.
\usepackage[autostyle, spanish=spanish]{csquotes} % Formatea bloques de citas.

% Bibliografía bajo norma UNE-ISO 690:2013.
\usepackage[
  backend=biber,     % "Backend" que genera la bibliografía.
  style=iso-numeric, % Usa el método de referencias numéricas.
  autolang=other,    % Soporta múltiples lenguajes en la bibliografía.
  sortlocale=es_ES,  % Ordena la bibliografía usando el lenguaje indicado.
  bibencoding=UTF8   % Codificación de caracteres de los archivos .bib.
]{biblatex}

% Indica el archivo usado por la memoria.
\addbibresource{bibliografia.bib}

\RequirePackage{booktabs}
\RequirePackage[table]{xcolor}
\RequirePackage{xtab}
\RequirePackage{multirow}

% Links
\usepackage[colorlinks]{hyperref}
\hypersetup{
	allcolors = {red}
}

% Ecuaciones
\usepackage{amsmath}

% Rutas de fichero / paquete
\newcommand{\ruta}[1]{{\sffamily #1}}

% Párrafos
\nonzeroparskip

% Imágenes
\usepackage{graphicx}
\newcommand{\imagen}[2]{
	\begin{figure}[!h]
		\centering
		\includegraphics[width=0.9\textwidth]{#1}
		\caption{#2}\label{fig:#1}
	\end{figure}
	\FloatBarrier
}

\newcommand{\imagenflotante}[2]{
	\begin{figure}%[!h]
		\centering
		\includegraphics[width=0.9\textwidth]{#1}
		\caption{#2}\label{fig:#1}
	\end{figure}
}

% El comando \figura nos permite insertar figuras cómodamente, y utilizando
% siempre el mismo formato. Los parámetros son:
% 1 -> Porcentaje del ancho de página que ocupará la figura (de 0 a 1)
% 2 --> Fichero de la imagen
% 3 --> Texto a pie de imagen
% 4 --> Etiqueta (label) para referencias
% 5 --> Opciones que queramos pasar a \includegraphics
% 6 --> Opciones de posicionamiento a pasar a \begin{figure}
\newcommand{\figuraConPosicion}[6]{%
  \setlength{\anchoFloat}{#1\textwidth}%
  \addtolength{\anchoFloat}{-4\fboxsep}%
  \setlength{\anchoFigura}{\anchoFloat}%
  \begin{figure}[#6]
    \begin{center}%
      \Ovalbox{%
        \begin{minipage}{\anchoFloat}%
          \begin{center}%
            \includegraphics[width=\anchoFigura,#5]{#2}%
            \caption{#3}%
            \label{#4}%
          \end{center}%
        \end{minipage}
      }%
    \end{center}%
  \end{figure}%
}

%
% Comando para incluir imágenes en formato apaisado (sin marco).
\newcommand{\figuraApaisadaSinMarco}[5]{%
  \begin{figure}%
    \begin{center}%
    \includegraphics[angle=90,height=#1\textheight,#5]{#2}%
    \caption{#3}%
    \label{#4}%
    \end{center}%
  \end{figure}%
}
% Para las tablas
\newcommand{\otoprule}{\midrule [\heavyrulewidth]}
%
% Nuevo comando para tablas pequeñas (menos de una página).
\newcommand{\tablaSmall}[5]{%
 \begin{table}
  \begin{center}
   \rowcolors {2}{gray!35}{}
   \begin{tabular}{#2}
    \toprule
    #4
    \otoprule
    #5
    \bottomrule
   \end{tabular}
   \caption{#1}
   \label{tabla:#3}
  \end{center}
 \end{table}
}

%
% Nuevo comando para tablas pequeñas (menos de una página).
\newcommand{\tablaSmallSinColores}[5]{%
 \begin{table}[H]
  \begin{center}
   \begin{tabular}{#2}
    \toprule
    #4
    \otoprule
    #5
    \bottomrule
   \end{tabular}
   \caption{#1}
   \label{tabla:#3}
  \end{center}
 \end{table}
}

\newcommand{\tablaApaisadaSmall}[5]{%
\begin{landscape}
  \begin{table}
   \begin{center}
    \rowcolors {2}{gray!35}{}
    \begin{tabular}{#2}
     \toprule
     #4
     \otoprule
     #5
     \bottomrule
    \end{tabular}
    \caption{#1}
    \label{tabla:#3}
   \end{center}
  \end{table}
\end{landscape}
}

%
% Nuevo comando para tablas grandes con cabecera y filas alternas coloreadas en gris.
\newcommand{\tabla}[6]{%
  \begin{center}
    \tablefirsthead{
      \toprule
      #5
      \otoprule
    }
    \tablehead{
      \multicolumn{#3}{l}{\small\sl continúa desde la página anterior}\\
      \toprule
      #5
      \otoprule
    }
    \tabletail{
      \hline
      \multicolumn{#3}{r}{\small\sl continúa en la página siguiente}\\
    }
    \tablelasttail{
      \hline
    }
    \bottomcaption{#1}
    \rowcolors {2}{gray!35}{}
    \begin{xtabular}{#2}
      #6
      \bottomrule
    \end{xtabular}
    \label{tabla:#4}
  \end{center}
}

%
% Nuevo comando para tablas grandes con cabecera.
\newcommand{\tablaSinColores}[6]{%
  \begin{center}
    \tablefirsthead{
      \toprule
      #5
      \otoprule
    }
    \tablehead{
      \multicolumn{#3}{l}{\small\sl continúa desde la página anterior}\\
      \toprule
      #5
      \otoprule
    }
    \tabletail{
      \hline
      \multicolumn{#3}{r}{\small\sl continúa en la página siguiente}\\
    }
    \tablelasttail{
      \hline
    }
    \bottomcaption{#1}
    \begin{xtabular}{#2}
      #6
      \bottomrule
    \end{xtabular}
    \label{tabla:#4}
  \end{center}
}

%
% Nuevo comando para tablas grandes sin cabecera.
\newcommand{\tablaSinCabecera}[5]{%
  \begin{center}
    \tablefirsthead{
      \toprule
    }
    \tablehead{
      \multicolumn{#3}{l}{\small\sl continúa desde la página anterior}\\
      \hline
    }
    \tabletail{
      \hline
      \multicolumn{#3}{r}{\small\sl continúa en la página siguiente}\\
    }
    \tablelasttail{
      \hline
    }
    \bottomcaption{#1}
  \begin{xtabular}{#2}
    #5
   \bottomrule
  \end{xtabular}
  \label{tabla:#4}
  \end{center}
}



\definecolor{cgoLight}{HTML}{EEEEEE}
\definecolor{cgoExtralight}{HTML}{FFFFFF}

%
% Nuevo comando para tablas grandes sin cabecera.
\newcommand{\tablaSinCabeceraConBandas}[5]{%
  \begin{center}
    \tablefirsthead{
      \toprule
    }
    \tablehead{
      \multicolumn{#3}{l}{\small\sl continúa desde la página anterior}\\
      \hline
    }
    \tabletail{
      \hline
      \multicolumn{#3}{r}{\small\sl continúa en la página siguiente}\\
    }
    \tablelasttail{
      \hline
    }
    \bottomcaption{#1}
    \rowcolors[]{1}{cgoExtralight}{cgoLight}

  \begin{xtabular}{#2}
    #5
   \bottomrule
  \end{xtabular}
  \label{tabla:#4}
  \end{center}
}

\graphicspath{ {./img/} }

% Capítulos
\chapterstyle{bianchi}
\newcommand{\capitulo}[2]{
	\setcounter{chapter}{#1}
	\setcounter{section}{0}
	\chapter*{#2}
	\addcontentsline{toc}{chapter}{#2}
	\markboth{#2}{#2}
}

% Apéndices
\renewcommand{\appendixname}{Apéndice}
\renewcommand*\cftappendixname{\appendixname}

\newcommand{\apendice}[1]{
	%\renewcommand{\thechapter}{A}
	\chapter{#1}
}

\renewcommand*\cftappendixname{\appendixname\ }

% Formato de portada
\makeatletter
\usepackage{xcolor}
\newcommand{\tutor}[1]{\def\@tutor{#1}}
\newcommand{\course}[1]{\def\@course{#1}}
\definecolor{cpardoBox}{HTML}{E6E6FF}
\def\maketitle{
  \null
  \thispagestyle{empty}
  % Cabecera ----------------
\noindent\includegraphics[width=\textwidth]{cabecera}\vspace{1cm}%
  \vfill
  % Título proyecto y escudo informática ----------------
  \colorbox{cpardoBox}{%
    \begin{minipage}{.8\textwidth}
      \vspace{.5cm}\Large
      \begin{center}
      \textbf{TFM del Máster Universitario en Ingeniería Informática}\vspace{.6cm}\\
      \textbf{\LARGE\@title{}}
      \end{center}
      \vspace{.2cm}
    \end{minipage}

  }%
  \hfill\begin{minipage}{.20\textwidth}
    \includegraphics[width=\textwidth]{escudoInfor}
  \end{minipage}
  \vfill
  % Datos de alumno, curso y tutores ------------------
  \begin{center}%
  {%
    \noindent\LARGE
    Presentado por \@author{}\\ 
    en Universidad de Burgos --- \@date{}\\
    Tutor: \@tutor{}\\
  }%
  \end{center}%
  \null
  \cleardoublepage
  }
\makeatother

\newcommand{\nombre}{RPC} %%% cambio de comando
     
% Comando para formatear palabras procedentes de otros
% lenguajes distintos del castellano.
% Estas palabras se pueden escribir en letra cursiva 
\newcommand{\extranjerismo}[1]{\textit{#1}}
% o con comillas si no se dispone de cursiva.
%\newcommand{\extranjerismo}[1]{"{#1}"}

% Comando para formatear los títulos de otras obras de creación.
\newcommand{\titulo}[1]{\textit{#1}}

% Datos de portada
\title{Comunicación TCP/IP con sistemas empotrados}
\author{\nombre}
\tutor{AMG}
\date{\today}

\begin{document}

\maketitle

\newpage\null\thispagestyle{empty}

%%%%%%%%%%%%%%%%%%%%%%%%%%%%%%%%%%%%%%%%%%%%%%%%%%%%%%%%%%%%%%%%%%%%%%%%%%%%%%%%%%%%%%%%
\thispagestyle{empty}


\noindent\includegraphics[width=\textwidth]{cabecera}\vspace{1cm}

\noindent D. tutor, profesor del Departamento de Ingeniería Electromecánica, Área de Ingeniería de Sistemas y Automática.

\noindent Expone:

\noindent Que el alumno D. \nombre, con DNI 12345678Z, ha realizado el Trabajo final de Máster Universitario en Ingeniería Informática titulado Comunicación TCP/IP con sistemas empotrados. 

\noindent Y que dicho trabajo ha sido realizado por el alumno bajo la dirección del que suscribe, en virtud de lo cual se autoriza su presentación y defensa.

\begin{center} %\large
En Burgos, {\large \today}
\end{center}

\vfill\vfill\vfill

% Author and supervisor
%\begin{minipage}{0.45\textwidth}
%\begin{flushleft} %\large
%Vº. Bº. del Tutor:\\[2cm]
%D. nombre tutor
%\end{flushleft}
%\end{minipage}
%\hfill
%\begin{minipage}{0.45\textwidth}
%\begin{flushleft} %\large
%Vº. Bº. del co-tutor:\\[2cm]
%D. nombre co-tutor
%\end{flushleft}
%\end{minipage}
%\hfill
%
%\vfill

% para casos con solo un tutor comentar lo anterior
% y descomentar lo siguiente
Vº. Bº. del Tutor:\\[2cm]
D. AMG

\newpage\null\thispagestyle{empty}

\frontmatter

% Abstract en castellano
\renewcommand*\abstractname{Resumen}
\begin{abstract}
  Las placas de desarrollo facilitan el estudio y desarrollo de sistemas 
  empotrados. Un sistema empotrado conectado a una red de comunicaciones de
  datos obtiene una nueva vía de interacción con otros sistemas, ya sean
  empotrados o convencionales. Un sistema empotrado conectado permite
  interactuar de forma remota con él, pudiendo realizar entre otras operaciones
  la consulta de sus sensores o la activación de sus actuadores.

  En este proyecto se muestra como conectar una placa de desarrollo FRDM-K64F
  a una red de área local usando el conjunto de protocolos TCP/IP.
  Aprovechando tanto el \extranjerismo{hardware} del que dispone la placa como
  del \extranjerismo{hardware} conectado a ella, se ejemplifica la interacción
  con algunos de sus dispositivos a través de una aplicación web. Como
  dispositivos configurados para su uso remoto se encuentran el led integrado
  en la placa, una pantalla LCD y una placa de expansión.
\end{abstract}

\renewcommand*\abstractname{Descriptores}
\begin{abstract}
  Sistemas embebidos, placa desarrollo, familia de protocolos de internet,
  aplicación web.
\end{abstract}

\clearpage

% Abstract en inglés
\renewcommand*\abstractname{Abstract}
\begin{abstract}
  Development boards facilitate the study and development of embedded systems.
  An embedded systems connected to a data network obtains a new way of 
  interaction with other systems, whether embedded or conventional. A connected
  embedded system allows to interact remotely with it, being able to carry out, 
  among other operations, the query of its sensors or the activation of its
  actuators.
  
  This project shows how to connect a FRDM-K64F development board to a local
  area network using the TCP / IP protocol suite. Taking advantage of both
  the hardware available on the board and the hardware connected to it, 
  the interaction with some of its devices is exemplified through a web
  application. As devices configured for remote use are the LED integrated in
  the board, an LCD screen and an expansion board.
\end{abstract}

\renewcommand*\abstractname{Keywords}
\begin{abstract}
  Embedded systems, development board, Internet protocol suite,
  web application.
\end{abstract}

\clearpage

% Índices
\tableofcontents

\clearpage

\listoffigures

\clearpage

\listoftables
\clearpage

\mainmatter
\capitulo{1}{Introducción}\label{ch:introduccion}
Los sistemas empotrados o embebidos (SE) son sistemas diseñados para realiza una
función específica y, por tanto, solo realizan una o unas pocas tareas
concretas. Este hecho les diferencia de otros sistemas de propósito general como
ordenadores o teléfonos móviles, que son capaces de realizar multitud de tares
de diferente naturaleza. 

Como el propio nombre de los SE indica, este tipo de sistemas se suele encontrar
integrado en otros sistemas eléctricos o mecánicos de mayor envergadura y que se
encargan de controlar. Por esta razón, se requiere que los SE cuenten con 
ciertas características: tamaño reducido, bajo consumo o bajo coste. Requisitos
que a su vez provocan nuevos atributos, uno de ellos es menor potencia de
cálculo en comparación con otro tipo de sistemas. Además, en función de las
condiciones ambientales donde se encuentre el SE es posible que necesite ser
dotado de protección adicional ante situaciones especiales de temperatura,
humedad, vibración, etc.

Ciertos SE controlan el funcionamiento de otros sistemas que requieren que sus
operaciones se realicen de manera determinista. Es decir, una instrucción dada
se tiene que realizar de manera inmediata o con un retardo mínimo y conocido de
antemano. Para dar soporte a operaciones en tiempo real, los SE cuentan con
Sistemas Operativos en Tiempo Real (RTOS) que se encargan de repartir el tiempo
de ejecución de cada tarea, asignándolo en función de la prioridad de cada
tarea.

A nivel de \extranjerismo{hardware}, un microcontrolador es quien se encarga de
ejecutar las instrucciones programadas. Cada microcontrolador cuenta con un
microprocesador, con memoria y con determinados periféricos de entrada y salida.
Para que el SE se pueda comunicar con sus periféricos o con otros dispositivos
se emplean los buses serie de datos. Algunos de los buses más utilizados son
Inter-Integrated Circuit (I\textsuperscript{2}C), Serial Peripheral Interface
(SPI), o Universal Serial Bus (USB).

Otro medio que emplean los SE para comunicarse son las redes de comunicaciones
de datos a las que se puede conectar tanto por cable como de manera
inalámbrica. Existen varios protocolos para transmitir datos en red, un conjunto
muy común y conocido es la familia de protocolos de internet, también
conocido como pila o conjunto de protocolos TCP/IP
(Transmission Control Protocol / Internet Protocol). Esta familia de protocolos
es ampliamente utilizada para comunicar equipos en redes de área local (LAN), en
redes de área local inalámbricas (WLAN) y en todo el mundo a través de Internet.

Con el uso de TCP/IP se abre un abanico de nuevas funcionalidades permitiendo
realizar funciones, apoyadas en otros protocolos de la pila TCP/IP, que antes
no eran posibles. Por ejemplo, usando el Hypertext Transfer Protocol (HTTP) se
pueden transferir páginas web, con File Transfer Protocol (FTP) transferir
archivos, con Simple Mail Transfer Protocol (SMTP) enviar correos electrónicos
o con Message Queuing Telemetry Transport (MQTT) se pueden enviar
mensajes bajo el patrón de mensajería de publicar-suscribir.

En este trabajo se muestra como crear un SE conectado usando TCP/IP en una placa
de desarrollo FRDM-K64F del fabricante NXP. Se parte de la configuración de los
componentes \extranjerismo{hardware} y \extranjerismo{software} necesarios, para
terminar demostrando la interacción, de manera remota, con el SE desde una
aplicación web. Desde dicha aplicación es posible enviar comandos para realizar
alguna de las funciones programadas en la placa.

\section{Estructura de la memoria}\label{sec:estructura}
La presente memoria se estructura de la siguiente manera:

\begin{itemize}
\item
  \textbf{Introducción:} descripción abreviada de los temas principales 
  abordados en el proyecto. La introducción está acompañada de la estructura
  que toma la memoria y un listado con el contenido adjunto a la misma.
\item
  \textbf{Objetivos del proyecto:} declaración de los objetivos generales,
  técnicos y personales que se pretenden conseguir con el desarrollo de este
  trabajo.
\item
  \textbf{Conceptos teóricos:} explicación de los conceptos teóricos más
  relevantes en la realización del proyecto. Se tratan tanto los sistemas
  empotrados como la familia de protocolos TCP/IP.
\item
  \textbf{Técnicas y herramientas:} descripción de las técnicas y las
  herramientas que han sido empleadas para el desarrollo del proyecto. También
  se muestran las alternativas valoradas y los motivos de escoger la opción
  seleccionada.
\item
  \textbf{Aspectos relevantes del desarrollo:} presentación de los aspectos o
  facetas que han tomado mayor relevancia durante la ejecución de trabajo.
\item
  \textbf{Trabajos relacionados:} estado de la técnica en la creación de
  sistemas que emplean sistemas empotrados conectados en red.
\item
  \textbf{Conclusiones y líneas de trabajo futuras:} conclusiones extraídas tras
  la realización del proyecto, así como nuevas líneas de trabajo sobre las que
  mejorar o ampliar lo presentando en este proyecto.
\end{itemize}

\section{Anexos a la memoria}\label{sec:anexos}
La memoria se presenta acompañada de los siguientes anexos:

\begin{itemize}
\item
  \textbf{Plan del proyecto software:} exposición de la
  planificación temporal y del estudio de viabilidad del proyecto, tanto la
  económica como la legal.
\item
  \textbf{Especificación de requisitos del software:}
  presentación del catálogo de requisitos, así como la descripción de cada uno
  de ellos.
\item
  \textbf{Especificación de diseño:} descripción de la fase de diseño, el
  diseño de datos, el diseño procedimental y el diseño arquitectónico, del
  \extranjerismo{software} del SE y de la aplicación web.
\item
  \textbf{Manual del programador:} explicación en detalle de aquellos aspectos
  que un programador debe conocer para trabajar con el código fuente del
  proyecto.
\item
  \textbf{Manual de usuario:} explicación para que un usuario interesado en el
  proyecto sea capaz de instalar, configurar y operar con el SE y con la
  aplicación web.
\end{itemize}

\clearpage

\section{Contenido adjunto}\label{sec:adjunto}
Se adjunta el siguiente contenido a la memoria y los anexos:

\begin{itemize}
\item
  \extranjerismo{Software} para la placa de desarrollo \cite{webpage:repo-se}.
\item
  Aplicación web para interacción remota \cite{webpage:repo-aw}.
\item	
  Documentación del \extranjerismo{software} \cite{webpage:repo-se-doc}.
\item	
  Documentación de la aplicación web \cite{webpage:repo-aw-doc}.
\item
  Repositorio en línea con el código del \extranjerismo{software} \cite{webpage:repo-se}.
\item
  Repositorio en línea con el código de la aplicación web \cite{webpage:repo-aw}.
\end{itemize}

\capitulo{2}{Objetivos del proyecto}\label{ch:objetivos}

En este capítulo se detallan los objetivos que se pretenden conseguir con la
ejecución del proyecto. Se diferencian tres tipos de objetivos, los generales
que dan causa al proyecto, los técnicos inherentes al tipo de proyecto realizado
y los personales que se desean conseguir \extranjerismo{motu proprio}.

\section{Objetivos generales}\label{sec:obj_generales}
\begin{itemize}
  \item Configurar un sistema empotrado que sea capaz de conectarse en red.
  \item Dotar al sistema empotrado de diversas funciones usando algunos de los
  periféricos de los que dispone.
  \item Demostrar la ejecución correcta de las funciones implementadas.
  \item Crear una interfaz web que permita interactuar con el sistema empotrado.
\end{itemize}

\section{Objetivos técnicos}\label{sec:obj_tecnicos}
\begin{itemize}
  \item Configurar el \extranjerismo{hardware} de una placa de desarrollo
  FRDM-K64F para poder conectarla a través de su puerto Ethernet a una LAN.
  \item Utilizar componentes \extranjerismo{hardware} de la placa como 
  son sus LED de colores para mostrar que es capaz de comunicarse via TCP/IP.
  \item Utilizar el bus serie de datos I\textsuperscript{2}C presente en la
  placa para extender su funcionalidad y que pueda mostrar mensajes en una
  pantalla.
  \item Emplear modulación por ancho de pulsos para regular la intensidad del
  brillo de unos LED presentes en una placa de expansión.
  \item A nivel de comunicaciones, usar la implementación ligera de 
  TCP/IP ``lwIP'' para manejar las transmisiones.
  \item Crear una aplicación web usando la tecnología JSF capaz de comunicarse
  con una placa conectada en red.
\end{itemize}

\section{Objetivos personales}\label{sec:obj_personales}
\begin{itemize}
  \item Ampliar los conocimientos y la experiencia en el desarrollo de
  \extranjerismo{software} para sistemas empotrados.
  \item Conocer el conjunto de herramientas de trabajo que proporciona el
  fabricante de la placa de desarrollo.
  \item Extender el entendimiento en la familia de protocolos TCP/IP.
  \item Revisar el proceso de desarrollo de una aplicación web.
  \item Emplear el sistema de control de versiones distribuido Git a través
  de la plataforma de desarrollo GitHub.
  \item Emplear la técnica de desarrollo ágil Scrum en las diferentes fases del
  desarrollo.
  \item Aprender a usar el sistema de composición de textos \LaTeX\ y utilizarlo
  para realizar la documentación del proyecto.
\end{itemize}

\capitulo{3}{Conceptos teóricos}
En este capítulo se sintetizan algunos de los aspectos tratados en este proyecto
para mejorar su compresión y entendimiento.

\section{Sistemas empotrados}
Parte importante del proyecto se centra en obtener un sistema empotrado capaz de
comunicarse usando los protocolos TCP/IP. A continuación se describen los
conceptos más relevantes en torno a los SE.

\subsection{Descripción}
Se puede considerar que un sistema empotrado es aquel cuyo
\extranjerismo{hardware} y \extranjerismo{software} se encuentran estrechamente
relacionados, está diseñado para cumplir con una función específica, se haya
integrado en un sistema mayor, no se espera que el usuario lo modifique y puede
trabajar sin interacción o con la mínima interacción humana necesaria.
\cite{jime13}

\imagen{refrigerator}{Ejemplo de integración de SE. \cite{nxp01}}

En la figura \ref{fig:refrigerator} se muestra un ejemplo de uso de varios SE
dentro de un sistema mayor, en este caso un frigorífico inteligente. El sistema
cuenta con varios componentes, la interfaz de usuario, la gestión del sistema,
el control del motor o la conectividad. Cada uno de los componentes se ayuda de
un SE para realizar su función asignada.

Los SE cuentan con el \extranjerismo{hardware} específico para la tarea a
realizar. La interfaz de usuario puede contar con pantallas o botones. La
gestión del sistema tiene acceso a sensores, el control de la corriente, la
iluminación o los ventiladores. El control del motor presenta componentes
eléctricos para la regulación del compresor. Un módulo de energía, un triodo
para corriente alterna o un módulo de corrección del factor de potencia son
componentes que pueden estar presentes en el control del motor. Asimismo, el
\extranjerismo{software} ejecutado en cada uno de estos SE varía según la
función a desempeñar realizando únicamente las tareas necesarias.

Por otra parte, también se puede advertir que los SE no están pensados para que
el usuario los modifique o programe, ni requieren que la interacción sea
constante por parte del usuario para su correcto funcionamiento.

Los SE y los sistemas que los emplean se encuentran fácilmente. Se hallan en 
sistemas de movilidad y transporte, automatización industrial, sector sanitario,
edificios inteligentes, redes de suministro inteligentes, investigación
científica, seguridad pública, supervisión de salud estructural, recuperación de
desastres, robótica, agricultura y ganadería, aplicaciones militares,
telecomunicaciones y electrónica de consumo. \cite{marw18}

\subsection{Características del \extranjerismo{hardware}}
Los SE disponen de componentes \extranjerismo{hardware} que siendo específicos
para la tarea a la que están destinados se pueden generalizar en procesador,
memoria y puertos de entrada y salida. El procesador se encarga de ejecutar las
instrucciones de los programas que manejan las entradas y las salidas del
sistema. Los programas ejecutados y los datos generados se almacenan en la
memoria. Y los puertos de entrada y salida, envían y reciben la señales con las
que trabaja el procesador.\cite{jime13}

También existen otros elementos que se encuentran a menudo en los SE:
\begin{itemize}
    \item Puertos de comunicación serie o paralelo
    \item Dispositivos de interfaz humana
    \item Sensores
    \item Actuadores
    \item Conversores analógica-digital (ADC)
    \item Conversores digital-analógica (DAC)
    \item Componentes de diagnóstico y redundancia
    \item Componentes de apoyo al sistema
    \item Otros subsistemas:
    \begin{itemize}
        \item Circuito integrado para aplicaciones específicas (ASIC)
        \item Matriz de puertas programables (FPGA)
    \end{itemize}
\end{itemize}

En la figura \ref{fig:diagr_bloques_hw} se presenta un diagrama de bloques
en el que se puede observar de forma general los componentes que forman un SE
y su funcionamiento. Las entradas son procesadas por el microcontrolador que a
su vez generará las salidas apropiadas. 
\imagen{diagr_bloques_hw}{Diagrama de bloques del HW de un SE}

\subsection{Características del \extranjerismo{software}}
En cuanto al \extranjerismo{software} presente en los sistemas embebidos, se
pueden diferenciar varios grupos. El primero de ellos los
\extranjerismo{drivers} o controladores encargados de la interacción directa
con el \extranjerismo{hardware} del SE.

Luego se encuentra el \extranjerismo{middleware}, que es aquel
\extranjerismo{software} que ocupa una posición entre el sistema operativo y los
programas. El término se usa con frecuencia en librerías de rutinas de
infraestructura que proporcionan servicios a los desarrolladores de los
programas. \cite{butt16}

Junto a lo anterior se puede hallar un Sistema operativo en tiempo real (RTOS).
En un sistema operativo de propósito general varias tareas se ejecutan de forma
aparentemente simultánea. De este modo se puede repartir el tiempo de ejecución
de manera equitativa entre usuarios, por ejemplo. En cambio, en un RTOS se prima
la ejecución de las tareas en un tiempo estrictamente limitado. Con un sistema
de prioridades se determina la importancia de las tareas y cuales necesitan ser
realizadas sin demora. \cite{amaz19}

Por último y funcionado sobre lo anterior se encuentra los programas necesarios
para el funcionamiento del SE. En caso de contar con un RTOS será este el que se
encargue de ejecutar una tarea u otra. Sino, de forma conocida como
\extranjerismo{bare-metal} se ejecutan las tareas de forma secuencial. Secuencia
que solo es alterada en caso surgir una interrupción.

\imagen{diagr_bloques_sw}{Digrama de bloques del SW de un SE}
En el diagrama de bloques de la figura \ref{fig:diagr_bloques_sw} se observan
como se ubican unos componentes sobre otros. Los componentes se comunican con
aquellos adyacentes. El RTOS es un componente opcional y que funciona de forma
paralela a \extranjerismo{drivers} y \extranjerismo{middleware}.

\capitulo{4}{Técnicas y herramientas}{\label{sec:tecyherra}}
Este capítulo presenta las técnicas metodológicas y las herramientas de
desarrollo usadas durante la realización del proyecto. En algunas de las
técnicas y herramientas existen varias alternativas a utilizar. En este capítulo
se comentan las características principales de cada alternativa y se expone el
porqué de la elección tomada.

\section{Técnicas metodológicas}{\label{sec:tecnicas}}

\subsubsection{Scrum}{\label{sec:scrum}}
Scrum es definido por dos de sus mayores promotores \textcite{schwaber2017}
como:
\begin{quotation}``un marco de trabajo en el que las personas pueden
  abordar problemas flexibles y complejos, mientras se entregan productos creativa y
  productivamente del mayor valor posible''
\end{quotation}.

Desde el punto de vista del desarrollo del \extranjerismo{software}, trabajar en
un producto en toda su extensión puede atraer problemas debidos, por ejemplo, a
cambios en los requisitos de los clientes. Para solventar este problema, Scrum
enuncia la realización del trabajo de forma ligera, iterativa e incremental,
permitiendo que el desarrollo pueda responder mejor ante imprevistos.

Scrum propone grupos de trabajo donde sus integrantes tienen diferentes roles y
responsabilidades. El grupo sigue un flujo de trabajo que gira entorno al
concepto de \extranjerismo{sprint}, definido como la unidad básica de
desarrollo. Antes de cada \extranjerismo{sprint} se planifica su duración
temporal, se definen la tareas a realizar en él y cuando ya está en marcha se
revisa diariamente para analizar su evolución. Una vez terminado se comprueban
los resultados obtenidos y se proporcionan los entregables generados.

\section{Herramientas \extranjerismo{hardware} de desarrollo}
  {\label{sec:herramientas}}

\subsubsection{Placa de desarrollo FRDM-K64F}{\label{sec:k64f}}
El uso de placas de desarrollo sirve como toma de contacto al microprocesador y
al resto del \extranjerismo{hardware} que componen un SE. En este proyecto la
placa utilizada es el modelo
\href{https://www.nxp.com/support/developer-resources/evaluation-and-development-boards/freedom-development-boards/mcu-boards/freedom-development-platform-for-kinetis-k64-k63-and-k24-mcus:FRDM-K64F}
{FRDM-K64F}.

\imagen{k64f}{Placa de desarrollo FRDM-K64F \cite{webpage:k64f}}

El microcontrolador (MCU) de esta placa es el \href{https://www.nxp.com/products/processors-and-microcontrollers/arm-based-processors-and-mcus/kinetis-cortex-m-mcus/k-seriesperformancem4/k6x-ethernet/kinetis-k64-120-mhz-256kb-sram-microcontrollers-mcus-based-on-arm-cortex-m4-core:K64_120}
{Kinetis K64}. Este MCU cuenta con el microprocesador \href{https://developer.arm.com/products/processors/cortex-m/cortex-m4}
{Cortex-M4} capaz de funcionar a 120 Mhz. Destaca la capacidad de funcionar en 
modos de energía ultra bajos, su prominente rendimiento y los dispositivos
integrados de conectividad y comunicación como los puertos General Purpose
Input/Output (GPIO), la interfaz Ethernet o los buses de datos serie
I\textsuperscript{2}C o SPI.

La placa K64F ofrece el \extranjerismo{hardware} necesario para poder aprovechar
las características que ofrece el MCU. Cuenta con los pines suficientes para dar 
cabida a los GPIO. La disposición de los conectores es compatible con la
utilizada en placas \href{https://www.arduino.cc/}{Arduino} permitiendo el uso
de multitud de placas de extensión ya existentes.

También presenta un puerto Ethernet que permite conectar la placa a una red
local y, a la postre, a Internet. También cuenta con \extranjerismo{hardware}
extra que posibilita nuevas opciones a la hora de investigar y desarrollar con
la placa. Por ejemplo, puertos USB, para alimentación, depuración y conexión, y
diodos emisor de luz (LED) de colores rojo, verde y azul (RGB).

\subsubsection{Placa de expansión Arduino Basic I/O}{\label{sec:basic-io}}
Como la placa K64F es compatible con las placas de expansión de Arduino se
utiliza la placa \href{https://web.archive.org/web/20160818213905/http://www.msebilbao.com/tienda/product_info.php?cPath=130&products_id=793&osCsid=f967e6ddeaaa2f19050972ff62295a08}
{Arduino Basic I/O} para poder utilizar \extranjerismo{hardware} adicional.

\imagen{basic_io}{Placa de expansión Arduino Basic I/O \cite{webpage:basio-io}}

En concreto, se utilizan sus 4 LED de colores que usados en combinación de la
técnica de modulación de señales pulse-width modulation (PWM) se puede regular
la intensidad del brillo de cada uno de los LED.

\subsubsection{Pantalla de cristal líquido (LCD)}{\label{sec:lcd}}
El disponer de una pantalla permite mostrar información en forma de texto a un
usuario del SE. En particular, esta pantalla puede mostrar 2 líneas de texto
de 16 caracteres cada una. Gracias a esta pantalla el usuario puede
conocer ciertos datos acerca del estado de inicialización de la placa, o puede
usarla para mostrar sus propios mensajes.

Cabe destacar que la pantalla usa una conexión de tipo paralela, necesitando
hasta 16 pines para la transmisión de datos, la alimentación, la iluminación y
el control de la transmisión. Como forma de simplificar su uso la pantalla
incorpora un módulo de comunicación I\textsuperscript{2}C, de tal manera que
solo es necesario el empleo de 4 pines.

\subsubsection{Accesorios extra}{\label{sec:extras}}
Para realizar el montaje de las placas y la pantalla se utilizan varios
accesorios extra.

Las placas de pruebas \cite{webpage:placa-pruebas} son unos tableros con
orificios que permiten conectar otras placas, componentes electrónicos,
conectores y cables. Su orificios están conectados eléctricamente, siguiendo un
patrón de bloques y líneas, permitiendo así la conexión entre los elementos del
sistema a montar. La rapidez y facilidad del montaje y desmontaje de los
elementos facilita el prototipado de sistemas antes de su fabricación.

Con el fin de conectar los pines de los componentes se usan varios cables
puente \cite{webpage:cable-puente}. Algunos conectados directamente a las
placas, otros conectados a las placas de pruebas. Para evitar confusiones los
cables son de colores, de tal manera que es posible distinguir la función que
realiza cada uno.

Por último, como el objetivo del proyecto es realizar un SE que utilice los
protocolos TCP/IP, es necesario conectar la placa K64F con un cable Ethernet
a un determinado equipamiento de acceso a la red, por ejemplo,
un \extranjerismo{switch}.

\capitulo{5}{Aspectos relevantes del desarrollo del proyecto}

Este apartado pretende recoger los aspectos más interesantes del desarrollo del proyecto, comentados por los autores del mismo.
Debe incluir desde la exposición del ciclo de vida utilizado, hasta los detalles de mayor relevancia de las fases de análisis, diseño e implementación.
Se busca que no sea una mera operación de copiar y pegar diagramas y extractos del código fuente, sino que realmente se justifiquen los caminos de solución que se han tomado, especialmente aquellos que no sean triviales.
Puede ser el lugar más adecuado para documentar los aspectos más interesantes del diseño y de la implementación, con un mayor hincapié en aspectos tales como el tipo de arquitectura elegido, los índices de las tablas de la base de datos, normalización y desnormalización, distribución en ficheros3, reglas de negocio dentro de las bases de datos (EDVHV GH GDWRV DFWLYDV), aspectos de desarrollo relacionados con el WWW...
Este apartado, debe convertirse en el resumen de la experiencia práctica del proyecto, y por sí mismo justifica que la memoria se convierta en un documento útil, fuente de referencia para los autores, los tutores y futuros alumnos.

\capitulo{6}{Trabajos relacionados}{\label{ch:relacionados}}
En el campo de los SE se han realizado sistemas complejos y de gran valor en 
numerosos ámbitos. En este sentido, \href{https://www.nxp.com/products/processors-and-microcontrollers/arm-based-processors-and-mcus/kinetis-cortex-m-mcus/k-seriesperformancem4/k6x-ethernet/kinetis-k64-120-mhz-256kb-sram-microcontrollers-mcus-based-on-arm-cortex-m4-core:K64_120}
{NXP} muestra varias aplicaciones de los SE en las que se pueden emplear el
mismo MCU, Kinetis~\textsuperscript{\tiny\textregistered} K64, existente en la
placa de desarrollo.

Los sistemas expuestos muestran características similares a las presentes en la
placa K64F y a las usadas en este proyecto. El uso de conectividad en red via
TCP/IP hace posible la comunicación con el SE a distancia. Y la utilización de
GPIO, I\textsuperscript{2}C o PWM permite operar con los sensores, actuadores u
otros componentes del SE.

A continuación se muestran cuatro sistemas pertenecientes a diferentes sectores
de aplicación. Los ejemplos pertenecen al sector industrial, sanitario y
comercial. En concreto, los sistemas operados por SE son los siguientes:

\begin{itemize}
  \item Vehículo no tripulado
  \item Sistema de telesalud
  \item Cama de hospital eléctrica
  \item Impresora de terminal punto de venta
\end{itemize}

La funcionalidad, desempeño y ámbito de uso de las aplicaciones resultan muy
diversas. Así pues, queda patente la versatilidad y el amplio espectro de
entornos en los que se pueden encontrar SE.

\subsection{Vehículo no tripulado}{\label{sec:uv}}
Un vehículo no tripulado es aquel vehículo capaz de funcionar sin un humano a
bordo. De forma general los vehículos no tripulados se clasifican en función del
medio en el que operan. Pueden ser terrestres, aéreos, acuáticos, submarinos,
espaciales... Independientemente del tipo que sean y de que estén operados por
control remoto o sean autónomos, los vehículos se sirven de SE para controlar
sus funciones.

\imagen{uv}{Diagrama de bloques de un vehículo no tripulado \cite{webpage:uv}}

Con LED controlados por GPIO se puede conocer el estado del vehículo. Para
dotar de conectividad al vehículo están disponibles interfaces cableadas e
inalámbricas.

El uso de PWM es una característica compartida con el proyecto. Mientras que 
en el proyecto se ha usado para regular la intensidad del brillo de unas
luces, en los vehículos no tripulados se emplea PWM para regular sus motores

Aunque no se ha incluido en el proyecto, la capacidad de usar tarjetas de
memoria, de dispositivos conectados a un puerto USB o la comunicación con
periféricos usando SPI está presente tanto en vehículos como en la placa de
desarrollo.

\subsection{Sistema de telesalud}{\label{sec:salud}}
El acceso a la asistencia sanitaria no siempre es factible o asequible. Vivir
en entornos rurales alejados de centros médicos, la falta de medios de
transporte, la pérdida de movilidad de los pacientes, o la falta de recursos
económicos son causas que pueden impedir dicha asistencia. Con el avance de la
tecnología y las telecomunicaciones ha surgido la posibilidad de abordar dicha
problemática con el uso de sistemas de telesalud.

Con el uso de sensores de diverso tipo se recogen datos sobre el estado de salud
del paciente. Su presión arterial, su ritmo cardíaco, su temperatura corporal,
etc. Los datos son recogido por el sistema, procesados y retransmitidos al
personal médico encargado de velar por la salud del paciente.

En la figura \ref{fig:telesalud} se muestran los componentes que integran el
sistema. Entorno al MCU se hallan presentes una interfaz infrarroja que recoge
los datos de los sensores y una pantalla que muestra información al paciente.
También se puede ver como el sistema usa PWM para generar sonidos.

Periódicamente el sistema envía los datos bien usando una conexión Ethernet o
bien usando una conexión USB a un ordenador conectado a Internet. De manera 
opcional se puede conectar el sistema a Internet de forma inalámbrica. Este
hecho serviría para aumentar la portabilidad del sistema de telesalud al
completo.

\imagen{telesalud}{Diagrama de bloques de un sistema de telesalud
\cite{webpage:telehealth}}

\subsection{Cama de hospital eléctrica}{\label{sec:cama}}
Siguiendo en el sector sanitario, se presenta otra muestra de aplicación de SE
en las camas de hospital eléctricas. El uso de SE permite ofrecer el 
máximo confort posible al paciente y mejorar atención sanitaria recibida.

Con una serie de elementos de regulación se puede acomodar la cama a las
necesidades del usuario. Para conseguir tal comodidad el SE cuenta con motores
encargados de regular la inclinación o altura de varias partes de la cama. 

Al integrar un monitor de constantes vitales o una bomba de infusión, 
dispostivos también creados con SE, se puede realizar una monitorización y 
atención completa del paciente. Y si además se conecta la cama a una red
inalámbrica se abre la posibilidad de monitorizar y atender remotamente al
paciente. Por ejemplo, desde un control de enfermería.

Para comunicarse telefónicamente con el enfermo se utiliza la conexión
Ethernet en conjunción con la tecnología voz sobre protocolo de internet (VoIP).

Por último, que la cama disponga de una pantalla LCD permite que conocer 
\extranjerismo{in situ} el estado de todo el sistema, paciente incluido.

\imagen{cama}{Diagrama de bloques de una cama de hospital eléctrica
\cite{webpage:bed}}

\subsection{Impresora de terminal punto de venta (TPV)}{\label{sec:tpv}}
Para terminar, se pasa al ámbito del sector minorista. Se muestra el sistema
utilizado en una impresora de un TPV. La impresora está compuesta por
dos SE. Mientras que uno de ellos se encarga de las entradas digitales, el
otro se ocupa del control de los motores y dispositivos de impresión.

Las opciones de conectividad permiten comunicarse con la impresora de varias
maneras. Con la conexión Ethernet se puede trabajar en red con la impresora.
Igualmente, se pueden usar las conexión USB o RS-485.

Para controlar el \extranjerismo{hardware} relativo al proceso de impresión
se cuenta con la ayuda de un segundo SE. Como el SE auxiliar recibe señales
analógicas requiere usar ADC, el resto de entradas y salidas se conectan
directamente a los pines GPIO. Finalmente, la comunicación con el SE principal
se realiza mediante I\textsuperscript{2}C o SPI.

Por otra parte, el SE principal cuenta con mayor número de conexiones al
recibir las entradas del usuario, los controles digitales, gestionar la
conectividad y manejar los sensores.

De nuevo están presentes características utilizadas a su vez en el proyecto:
conectividad en red a través de un puerto Ethernet, conexión con periféricos
usando el bus de datos serie I\textsuperscript{2}C y comunicaciones diversas via
GPIO.

\imagen{tpv}{Diagrama de bloques de una impresora de TPV\cite{webpage:tpv}}

\capitulo{7}{Conclusiones y Líneas de trabajo futuras}

Todo proyecto debe incluir las conclusiones que se derivan de su desarrollo. Éstas pueden ser de diferente índole, dependiendo de la tipología del proyecto, pero normalmente van a estar presentes un conjunto de conclusiones relacionadas con los resultados del proyecto y un conjunto de conclusiones técnicas. 
Además, resulta muy útil realizar un informe crítico indicando cómo se puede mejorar el proyecto, o cómo se puede continuar trabajando en la línea del proyecto realizado. 


% Genera la bibliografía.
\printbibliography

\end{document}
